\section{Audit et identification des vulnérabilités}\label{sec:audit}

L'audit est limité aux adresses de la maquette.
Il commence par l'identification des hôtes et des services accessibles depuis Kali, puis par l'observation contrôlée des fonctions de DVWA\@.
Les commandes, dates et sorties utiles doivent être conservées dans le répertoire de preuves du groupe afin que les résultats soient reproductibles.
La reconnaissance réseau doit notamment confirmer que seul le service HTTP de DVWA est accessible depuis le réseau \texttt{attack}, et que la base MariaDB n'est pas directement atteignable.
Cette vérification contrôle aussi l'efficacité de la segmentation retenue.
Section~\ref{sec:maquette} présente les mécanismes mis en place.

\begin{table}[H]
    \centering
    \caption{Synthèse des surfaces d'attaque à vérifier.}\label{tab:audit}
    \begin{tabular}{@{}p{0.24\textwidth}p{0.27\textwidth}p{0.36\textwidth}@{}}
        \toprule
        Élément & Risque étudié & Validation attendue \\
        \midrule
        Interface web DVWA & Injection SQL & Confirmer le comportement vulnérable avec le niveau de sécurité retenu. \\
        Authentification DVWA & Force brute & Identifier la présence ou l'absence de mécanismes de limitation. \\
        Service MariaDB & Exposition réseau & Vérifier l'absence d'accès direct depuis Kali. \\
        Passerelle & Contournement du filtrage & Vérifier que seuls les flux HTTP prévus transitent. \\
        \bottomrule
    \end{tabular}
\end{table}

Les vulnérabilités de DVWA sont intentionnelles et documentées par le projet~\cite{dvwa}.
Elles ne dispensent pas de décrire précisément la configuration utilisée : niveau de sécurité DVWA, comptes de test, version des images et date des essais.
Ces paramètres conditionnent la reproductibilité et l'interprétation des résultats.
